% Options for packages loaded elsewhere
\PassOptionsToPackage{unicode}{hyperref}
\PassOptionsToPackage{hyphens}{url}
\documentclass[
]{article}
\usepackage{xcolor}
\usepackage[margin=1in]{geometry}
\usepackage{amsmath,amssymb}
\setcounter{secnumdepth}{-\maxdimen} % remove section numbering
\usepackage{iftex}
\ifPDFTeX
  \usepackage[T1]{fontenc}
  \usepackage[utf8]{inputenc}
  \usepackage{textcomp} % provide euro and other symbols
\else % if luatex or xetex
  \usepackage{unicode-math} % this also loads fontspec
  \defaultfontfeatures{Scale=MatchLowercase}
  \defaultfontfeatures[\rmfamily]{Ligatures=TeX,Scale=1}
\fi
\usepackage{lmodern}
\ifPDFTeX\else
  % xetex/luatex font selection
\fi
% Use upquote if available, for straight quotes in verbatim environments
\IfFileExists{upquote.sty}{\usepackage{upquote}}{}
\IfFileExists{microtype.sty}{% use microtype if available
  \usepackage[]{microtype}
  \UseMicrotypeSet[protrusion]{basicmath} % disable protrusion for tt fonts
}{}
\makeatletter
\@ifundefined{KOMAClassName}{% if non-KOMA class
  \IfFileExists{parskip.sty}{%
    \usepackage{parskip}
  }{% else
    \setlength{\parindent}{0pt}
    \setlength{\parskip}{6pt plus 2pt minus 1pt}}
}{% if KOMA class
  \KOMAoptions{parskip=half}}
\makeatother
\usepackage{color}
\usepackage{fancyvrb}
\newcommand{\VerbBar}{|}
\newcommand{\VERB}{\Verb[commandchars=\\\{\}]}
\DefineVerbatimEnvironment{Highlighting}{Verbatim}{commandchars=\\\{\}}
% Add ',fontsize=\small' for more characters per line
\usepackage{framed}
\definecolor{shadecolor}{RGB}{248,248,248}
\newenvironment{Shaded}{\begin{snugshade}}{\end{snugshade}}
\newcommand{\AlertTok}[1]{\textcolor[rgb]{0.94,0.16,0.16}{#1}}
\newcommand{\AnnotationTok}[1]{\textcolor[rgb]{0.56,0.35,0.01}{\textbf{\textit{#1}}}}
\newcommand{\AttributeTok}[1]{\textcolor[rgb]{0.13,0.29,0.53}{#1}}
\newcommand{\BaseNTok}[1]{\textcolor[rgb]{0.00,0.00,0.81}{#1}}
\newcommand{\BuiltInTok}[1]{#1}
\newcommand{\CharTok}[1]{\textcolor[rgb]{0.31,0.60,0.02}{#1}}
\newcommand{\CommentTok}[1]{\textcolor[rgb]{0.56,0.35,0.01}{\textit{#1}}}
\newcommand{\CommentVarTok}[1]{\textcolor[rgb]{0.56,0.35,0.01}{\textbf{\textit{#1}}}}
\newcommand{\ConstantTok}[1]{\textcolor[rgb]{0.56,0.35,0.01}{#1}}
\newcommand{\ControlFlowTok}[1]{\textcolor[rgb]{0.13,0.29,0.53}{\textbf{#1}}}
\newcommand{\DataTypeTok}[1]{\textcolor[rgb]{0.13,0.29,0.53}{#1}}
\newcommand{\DecValTok}[1]{\textcolor[rgb]{0.00,0.00,0.81}{#1}}
\newcommand{\DocumentationTok}[1]{\textcolor[rgb]{0.56,0.35,0.01}{\textbf{\textit{#1}}}}
\newcommand{\ErrorTok}[1]{\textcolor[rgb]{0.64,0.00,0.00}{\textbf{#1}}}
\newcommand{\ExtensionTok}[1]{#1}
\newcommand{\FloatTok}[1]{\textcolor[rgb]{0.00,0.00,0.81}{#1}}
\newcommand{\FunctionTok}[1]{\textcolor[rgb]{0.13,0.29,0.53}{\textbf{#1}}}
\newcommand{\ImportTok}[1]{#1}
\newcommand{\InformationTok}[1]{\textcolor[rgb]{0.56,0.35,0.01}{\textbf{\textit{#1}}}}
\newcommand{\KeywordTok}[1]{\textcolor[rgb]{0.13,0.29,0.53}{\textbf{#1}}}
\newcommand{\NormalTok}[1]{#1}
\newcommand{\OperatorTok}[1]{\textcolor[rgb]{0.81,0.36,0.00}{\textbf{#1}}}
\newcommand{\OtherTok}[1]{\textcolor[rgb]{0.56,0.35,0.01}{#1}}
\newcommand{\PreprocessorTok}[1]{\textcolor[rgb]{0.56,0.35,0.01}{\textit{#1}}}
\newcommand{\RegionMarkerTok}[1]{#1}
\newcommand{\SpecialCharTok}[1]{\textcolor[rgb]{0.81,0.36,0.00}{\textbf{#1}}}
\newcommand{\SpecialStringTok}[1]{\textcolor[rgb]{0.31,0.60,0.02}{#1}}
\newcommand{\StringTok}[1]{\textcolor[rgb]{0.31,0.60,0.02}{#1}}
\newcommand{\VariableTok}[1]{\textcolor[rgb]{0.00,0.00,0.00}{#1}}
\newcommand{\VerbatimStringTok}[1]{\textcolor[rgb]{0.31,0.60,0.02}{#1}}
\newcommand{\WarningTok}[1]{\textcolor[rgb]{0.56,0.35,0.01}{\textbf{\textit{#1}}}}
\usepackage{longtable,booktabs,array}
\usepackage{calc} % for calculating minipage widths
% Correct order of tables after \paragraph or \subparagraph
\usepackage{etoolbox}
\makeatletter
\patchcmd\longtable{\par}{\if@noskipsec\mbox{}\fi\par}{}{}
\makeatother
% Allow footnotes in longtable head/foot
\IfFileExists{footnotehyper.sty}{\usepackage{footnotehyper}}{\usepackage{footnote}}
\makesavenoteenv{longtable}
\usepackage{graphicx}
\makeatletter
\newsavebox\pandoc@box
\newcommand*\pandocbounded[1]{% scales image to fit in text height/width
  \sbox\pandoc@box{#1}%
  \Gscale@div\@tempa{\textheight}{\dimexpr\ht\pandoc@box+\dp\pandoc@box\relax}%
  \Gscale@div\@tempb{\linewidth}{\wd\pandoc@box}%
  \ifdim\@tempb\p@<\@tempa\p@\let\@tempa\@tempb\fi% select the smaller of both
  \ifdim\@tempa\p@<\p@\scalebox{\@tempa}{\usebox\pandoc@box}%
  \else\usebox{\pandoc@box}%
  \fi%
}
% Set default figure placement to htbp
\def\fps@figure{htbp}
\makeatother
\setlength{\emergencystretch}{3em} % prevent overfull lines
\providecommand{\tightlist}{%
  \setlength{\itemsep}{0pt}\setlength{\parskip}{0pt}}
\usepackage{bookmark}
\IfFileExists{xurl.sty}{\usepackage{xurl}}{} % add URL line breaks if available
\urlstyle{same}
\hypersetup{
  pdftitle={2024 MCM Problem A --- Resource Availability and Sex Ratios},
  pdfauthor={Generated by ChatGPT},
  hidelinks,
  pdfcreator={LaTeX via pandoc}}

\title{2024 MCM Problem A --- Resource Availability and Sex Ratios}
\author{Generated by ChatGPT}
\date{November 13, 2025}

\begin{document}
\maketitle

\section{One-Page Summary}\label{one-page-summary}

Question. Explore how variable sex ratios in sea lampreys, driven by
larval growth rates which depend on food availability, affect lamprey
population dynamics and the larger ecosystem. We investigate advantages
and disadvantages of adaptive sex ratios, impacts on stability, and
secondary effects on hosts and parasites.

Approach. We build a stage-structured dynamical model with resource
dynamics and a sex-determination function m(R) mapping resource
availability to fraction male. The core model couples: host fish (F),
lamprey population stages (eggs E, larvae L, adults A split by sex M and
W), and resource pool (R). We analyze equilibria, perform stability and
sensitivity analyses, and run numerical simulations under low and high
resource regimes.

Key results (summary).

\begin{itemize}
\tightlist
\item
  Adaptive sex ratio that produces more males under low food can protect
  the lamprey population by reducing reproductive output when resources
  are scarce, acting as a bet-hedging strategy and lowering the risk of
  overshooting the resource base.
\item
  Variable sex ratios change ecosystem stability: in some parameter
  regions they stabilize host-lamprey dynamics (damp oscillations) but
  can destabilize when reproduction responds nonlinearly to resources.
\item
  Secondary effects: variable sex ratios can indirectly benefit
  parasites (and other species) by smoothing lamprey-driven fish
  mortality peaks, changing predation pressure timing.
\end{itemize}

Conclusions. Adaptive sex ratio variation is a plausible evolutionary
response to variable resources; it mitigates boom-bust dynamics for
lampreys and alters ecosystem resilience. Management implications:
controlling resource availability (for example, habitat alteration) can
shift sex ratios and thus population trajectories.

\section{Model development}\label{model-development}

\subsection{Model variables}\label{model-variables}

R(t): dimensionless resource availability (normalized, 0 to 1). F(t):
density of host fish (food source/hosts for adult lampreys). E(t): egg
stock. L(t): larvae. M(t): adult males. W(t): adult females. A(t) = M(t)
+ W(t): total adult lampreys.

\subsection{Resource-dependent sex
ratio}\label{resource-dependent-sex-ratio}

We define a simple linear mapping from resource level R in {[}0,1{]} to
male fraction m(R) matching the observed endpoints:

m(R) = 0.78 - 0.22 * R

Thus m(0) = 0.78, m(1) = 0.56. Female fraction is 1 - m(R).

\subsection{ODE system
(stage-structured)}\label{ode-system-stage-structured}

The model equations are presented here (derivatives):

dR/dt = S - dR * R - cR * L - cA * A

dF/dt = rF * F * (1 - F/K) - beta * A * F / (F + H)

dE/dt = b * W - sigmaE * E - muE * E

dL/dt = sigmaE * E - sigmaL * L - muL * L

dM/dt = m(R) * sigmaL * L - muA * M

dW/dt = (1 - m(R)) * sigmaL * L - muA * W

Explanation of parameters

S: resource supply rate (controls low/high regimes)

dR: resource decay/loss

cR, cA: consumption coefficients by larvae and adults

rF: intrinsic growth rate of host fish, K carrying capacity

beta: mortality rate of fish due to lamprey parasitism (functional
response with half-saturation H)

b: per-female effective reproductive contribution (eggs that contribute
to E)

sigmaE, sigmaL: progression rates from eggs to larvae and larvae to
adults

muE, muL, muA: stage-specific mortality rates

\section{Numerical simulations}\label{numerical-simulations}

We simulate under two resource-supply regimes: low and high. The R code
below implements the ODE system using deSolve's ode solver.

\begin{Shaded}
\begin{Highlighting}[]
\CommentTok{\# Sex ratio function}
\NormalTok{mR }\OtherTok{\textless{}{-}} \ControlFlowTok{function}\NormalTok{(R) \{}
  \FunctionTok{pmax}\NormalTok{(}\DecValTok{0}\NormalTok{, }\FunctionTok{pmin}\NormalTok{(}\DecValTok{1}\NormalTok{, }\FloatTok{0.78} \SpecialCharTok{{-}} \FloatTok{0.22} \SpecialCharTok{*}\NormalTok{ R))}
\NormalTok{\}}

\CommentTok{\# Derivative function for deSolve}
\NormalTok{lamprey\_model }\OtherTok{\textless{}{-}} \ControlFlowTok{function}\NormalTok{(t, state, params) \{}
  \FunctionTok{with}\NormalTok{(}\FunctionTok{as.list}\NormalTok{(}\FunctionTok{c}\NormalTok{(state, params)), \{}
\NormalTok{    A }\OtherTok{\textless{}{-}}\NormalTok{ M }\SpecialCharTok{+}\NormalTok{ W}
\NormalTok{    dR }\OtherTok{\textless{}{-}}\NormalTok{ S }\SpecialCharTok{{-}}\NormalTok{ dR\_coef }\SpecialCharTok{*}\NormalTok{ R }\SpecialCharTok{{-}}\NormalTok{ cR }\SpecialCharTok{*}\NormalTok{ L }\SpecialCharTok{{-}}\NormalTok{ cA }\SpecialCharTok{*}\NormalTok{ A}
\NormalTok{    dF }\OtherTok{\textless{}{-}}\NormalTok{ rF }\SpecialCharTok{*}\NormalTok{ F }\SpecialCharTok{*}\NormalTok{ (}\DecValTok{1} \SpecialCharTok{{-}}\NormalTok{ F }\SpecialCharTok{/}\NormalTok{ K) }\SpecialCharTok{{-}}\NormalTok{ beta }\SpecialCharTok{*}\NormalTok{ A }\SpecialCharTok{*}\NormalTok{ (F }\SpecialCharTok{/}\NormalTok{ (F }\SpecialCharTok{+}\NormalTok{ H))}
\NormalTok{    dE }\OtherTok{\textless{}{-}}\NormalTok{ b }\SpecialCharTok{*}\NormalTok{ W }\SpecialCharTok{{-}}\NormalTok{ sigmaE }\SpecialCharTok{*}\NormalTok{ E }\SpecialCharTok{{-}}\NormalTok{ muE }\SpecialCharTok{*}\NormalTok{ E}
\NormalTok{    dL }\OtherTok{\textless{}{-}}\NormalTok{ sigmaE }\SpecialCharTok{*}\NormalTok{ E }\SpecialCharTok{{-}}\NormalTok{ sigmaL }\SpecialCharTok{*}\NormalTok{ L }\SpecialCharTok{{-}}\NormalTok{ muL }\SpecialCharTok{*}\NormalTok{ L}
\NormalTok{    mr }\OtherTok{\textless{}{-}} \FunctionTok{mR}\NormalTok{(R)}
\NormalTok{    dM }\OtherTok{\textless{}{-}}\NormalTok{ mr }\SpecialCharTok{*}\NormalTok{ sigmaL }\SpecialCharTok{*}\NormalTok{ L }\SpecialCharTok{{-}}\NormalTok{ muA }\SpecialCharTok{*}\NormalTok{ M}
\NormalTok{    dW }\OtherTok{\textless{}{-}}\NormalTok{ (}\DecValTok{1} \SpecialCharTok{{-}}\NormalTok{ mr) }\SpecialCharTok{*}\NormalTok{ sigmaL }\SpecialCharTok{*}\NormalTok{ L }\SpecialCharTok{{-}}\NormalTok{ muA }\SpecialCharTok{*}\NormalTok{ W}
    \FunctionTok{list}\NormalTok{(}\FunctionTok{c}\NormalTok{(dR, dF, dE, dL, dM, dW))}
\NormalTok{  \})}
\NormalTok{\}}

\CommentTok{\# Parameters and initial state}
\NormalTok{common }\OtherTok{\textless{}{-}} \FunctionTok{list}\NormalTok{(}\AttributeTok{dR\_coef =} \FloatTok{0.2}\NormalTok{, }\AttributeTok{cR =} \FloatTok{0.01}\NormalTok{, }\AttributeTok{cA =} \FloatTok{0.02}\NormalTok{, }\AttributeTok{rF =} \FloatTok{0.5}\NormalTok{, }\AttributeTok{K =} \DecValTok{100}\NormalTok{,}
               \AttributeTok{b =} \FloatTok{0.5}\NormalTok{, }\AttributeTok{beta =} \FloatTok{0.02}\NormalTok{, }\AttributeTok{H =} \DecValTok{10}\NormalTok{, }\AttributeTok{sigmaE =} \FloatTok{0.2}\NormalTok{, }\AttributeTok{sigmaL =} \FloatTok{0.1}\NormalTok{,}
               \AttributeTok{muE =} \FloatTok{0.1}\NormalTok{, }\AttributeTok{muL =} \FloatTok{0.05}\NormalTok{, }\AttributeTok{muA =} \FloatTok{0.1}\NormalTok{)}
\NormalTok{params\_low }\OtherTok{\textless{}{-}} \FunctionTok{c}\NormalTok{(}\AttributeTok{S =} \FloatTok{0.5}\NormalTok{, common)}
\NormalTok{params\_high }\OtherTok{\textless{}{-}} \FunctionTok{c}\NormalTok{(}\AttributeTok{S =} \FloatTok{1.5}\NormalTok{, common)}

\NormalTok{state0 }\OtherTok{\textless{}{-}} \FunctionTok{c}\NormalTok{(}\AttributeTok{R =} \FloatTok{0.5}\NormalTok{, }\AttributeTok{F =} \DecValTok{50}\NormalTok{, }\AttributeTok{E =} \DecValTok{10}\NormalTok{, }\AttributeTok{L =} \DecValTok{20}\NormalTok{, }\AttributeTok{M =} \DecValTok{5}\NormalTok{, }\AttributeTok{W =} \DecValTok{5}\NormalTok{)}

\NormalTok{times }\OtherTok{\textless{}{-}} \FunctionTok{seq}\NormalTok{(}\DecValTok{0}\NormalTok{, }\DecValTok{400}\NormalTok{, }\AttributeTok{by =} \FloatTok{0.1}\NormalTok{)}

\NormalTok{out\_low }\OtherTok{\textless{}{-}} \FunctionTok{ode}\NormalTok{(}\AttributeTok{y =}\NormalTok{ state0, }\AttributeTok{times =}\NormalTok{ times, }\AttributeTok{func =}\NormalTok{ lamprey\_model, }\AttributeTok{parms =}\NormalTok{ params\_low)}
\NormalTok{out\_high }\OtherTok{\textless{}{-}} \FunctionTok{ode}\NormalTok{(}\AttributeTok{y =}\NormalTok{ state0, }\AttributeTok{times =}\NormalTok{ times, }\AttributeTok{func =}\NormalTok{ lamprey\_model, }\AttributeTok{parms =}\NormalTok{ params\_high)}

\NormalTok{out\_low }\OtherTok{\textless{}{-}} \FunctionTok{as.data.frame}\NormalTok{(out\_low)}
\NormalTok{out\_high }\OtherTok{\textless{}{-}} \FunctionTok{as.data.frame}\NormalTok{(out\_high)}

\CommentTok{\# compute total adults and male fraction}
\NormalTok{out\_low}\SpecialCharTok{$}\NormalTok{A }\OtherTok{\textless{}{-}}\NormalTok{ out\_low}\SpecialCharTok{$}\NormalTok{M }\SpecialCharTok{+}\NormalTok{ out\_low}\SpecialCharTok{$}\NormalTok{W}
\NormalTok{out\_low}\SpecialCharTok{$}\NormalTok{m }\OtherTok{\textless{}{-}} \FloatTok{0.78} \SpecialCharTok{{-}} \FloatTok{0.22} \SpecialCharTok{*}\NormalTok{ out\_low}\SpecialCharTok{$}\NormalTok{R}
\NormalTok{out\_high}\SpecialCharTok{$}\NormalTok{A }\OtherTok{\textless{}{-}}\NormalTok{ out\_high}\SpecialCharTok{$}\NormalTok{M }\SpecialCharTok{+}\NormalTok{ out\_high}\SpecialCharTok{$}\NormalTok{W}
\NormalTok{out\_high}\SpecialCharTok{$}\NormalTok{m }\OtherTok{\textless{}{-}} \FloatTok{0.78} \SpecialCharTok{{-}} \FloatTok{0.22} \SpecialCharTok{*}\NormalTok{ out\_high}\SpecialCharTok{$}\NormalTok{R}

\CommentTok{\# Save plots}
\NormalTok{p1 }\OtherTok{\textless{}{-}} \FunctionTok{ggplot}\NormalTok{(out\_low, }\FunctionTok{aes}\NormalTok{(}\AttributeTok{x =}\NormalTok{ time)) }\SpecialCharTok{+}
  \FunctionTok{geom\_line}\NormalTok{(}\FunctionTok{aes}\NormalTok{(}\AttributeTok{y =}\NormalTok{ F, }\AttributeTok{colour =} \StringTok{\textquotesingle{}Fish\textquotesingle{}}\NormalTok{)) }\SpecialCharTok{+}
  \FunctionTok{geom\_line}\NormalTok{(}\FunctionTok{aes}\NormalTok{(}\AttributeTok{y =}\NormalTok{ A, }\AttributeTok{colour =} \StringTok{\textquotesingle{}Lampreys\textquotesingle{}}\NormalTok{)) }\SpecialCharTok{+}
  \FunctionTok{labs}\NormalTok{(}\AttributeTok{title =} \StringTok{\textquotesingle{}Low{-}resource scenario\textquotesingle{}}\NormalTok{, }\AttributeTok{y =} \StringTok{\textquotesingle{}Population\textquotesingle{}}\NormalTok{) }\SpecialCharTok{+}
  \FunctionTok{scale\_color\_manual}\NormalTok{(}\AttributeTok{values =} \FunctionTok{c}\NormalTok{(}\StringTok{\textquotesingle{}Fish\textquotesingle{}} \OtherTok{=} \StringTok{\textquotesingle{}black\textquotesingle{}}\NormalTok{, }\StringTok{\textquotesingle{}Lampreys\textquotesingle{}} \OtherTok{=} \StringTok{\textquotesingle{}grey\textquotesingle{}}\NormalTok{))}

\NormalTok{p2 }\OtherTok{\textless{}{-}} \FunctionTok{ggplot}\NormalTok{(out\_high, }\FunctionTok{aes}\NormalTok{(}\AttributeTok{x =}\NormalTok{ time)) }\SpecialCharTok{+}
  \FunctionTok{geom\_line}\NormalTok{(}\FunctionTok{aes}\NormalTok{(}\AttributeTok{y =}\NormalTok{ F, }\AttributeTok{colour =} \StringTok{\textquotesingle{}Fish\textquotesingle{}}\NormalTok{)) }\SpecialCharTok{+}
  \FunctionTok{geom\_line}\NormalTok{(}\FunctionTok{aes}\NormalTok{(}\AttributeTok{y =}\NormalTok{ A, }\AttributeTok{colour =} \StringTok{\textquotesingle{}Lampreys\textquotesingle{}}\NormalTok{)) }\SpecialCharTok{+}
  \FunctionTok{labs}\NormalTok{(}\AttributeTok{title =} \StringTok{\textquotesingle{}High{-}resource scenario\textquotesingle{}}\NormalTok{, }\AttributeTok{y =} \StringTok{\textquotesingle{}Population\textquotesingle{}}\NormalTok{) }\SpecialCharTok{+}
  \FunctionTok{scale\_color\_manual}\NormalTok{(}\AttributeTok{values =} \FunctionTok{c}\NormalTok{(}\StringTok{\textquotesingle{}Fish\textquotesingle{}} \OtherTok{=} \StringTok{\textquotesingle{}black\textquotesingle{}}\NormalTok{, }\StringTok{\textquotesingle{}Lampreys\textquotesingle{}} \OtherTok{=} \StringTok{\textquotesingle{}grey\textquotesingle{}}\NormalTok{))}

\NormalTok{p3 }\OtherTok{\textless{}{-}} \FunctionTok{ggplot}\NormalTok{() }\SpecialCharTok{+}
  \FunctionTok{geom\_line}\NormalTok{(}\AttributeTok{data =}\NormalTok{ out\_low, }\FunctionTok{aes}\NormalTok{(}\AttributeTok{x =}\NormalTok{ time, }\AttributeTok{y =}\NormalTok{ m, }\AttributeTok{colour =} \StringTok{\textquotesingle{}Low\textquotesingle{}}\NormalTok{)) }\SpecialCharTok{+}
  \FunctionTok{geom\_line}\NormalTok{(}\AttributeTok{data =}\NormalTok{ out\_high, }\FunctionTok{aes}\NormalTok{(}\AttributeTok{x =}\NormalTok{ time, }\AttributeTok{y =}\NormalTok{ m, }\AttributeTok{colour =} \StringTok{\textquotesingle{}High\textquotesingle{}}\NormalTok{)) }\SpecialCharTok{+}
  \FunctionTok{labs}\NormalTok{(}\AttributeTok{title =} \StringTok{\textquotesingle{}Male fraction over time\textquotesingle{}}\NormalTok{, }\AttributeTok{y =} \StringTok{\textquotesingle{}Male fraction m(R)\textquotesingle{}}\NormalTok{) }\SpecialCharTok{+}
  \FunctionTok{scale\_color\_manual}\NormalTok{(}\AttributeTok{values =} \FunctionTok{c}\NormalTok{(}\StringTok{\textquotesingle{}Low\textquotesingle{}} \OtherTok{=} \StringTok{\textquotesingle{}black\textquotesingle{}}\NormalTok{, }\StringTok{\textquotesingle{}High\textquotesingle{}} \OtherTok{=} \StringTok{\textquotesingle{}grey\textquotesingle{}}\NormalTok{))}
\end{Highlighting}
\end{Shaded}

\begin{Shaded}
\begin{Highlighting}[]
\NormalTok{p1}
\end{Highlighting}
\end{Shaded}

\pandocbounded{\includegraphics[keepaspectratio]{mcm_2024_problem_a_files/figure-latex/unnamed-chunk-2-1.pdf}}

\begin{Shaded}
\begin{Highlighting}[]
\NormalTok{p2}
\end{Highlighting}
\end{Shaded}

\pandocbounded{\includegraphics[keepaspectratio]{mcm_2024_problem_a_files/figure-latex/unnamed-chunk-3-1.pdf}}

\begin{Shaded}
\begin{Highlighting}[]
\NormalTok{p3}
\end{Highlighting}
\end{Shaded}

\pandocbounded{\includegraphics[keepaspectratio]{mcm_2024_problem_a_files/figure-latex/unnamed-chunk-4-1.pdf}}

\begin{Shaded}
\begin{Highlighting}[]
\CommentTok{\# Final summary table}
\NormalTok{final\_low }\OtherTok{\textless{}{-}} \FunctionTok{tail}\NormalTok{(out\_low, }\DecValTok{1}\NormalTok{)}
\NormalTok{final\_high }\OtherTok{\textless{}{-}} \FunctionTok{tail}\NormalTok{(out\_high, }\DecValTok{1}\NormalTok{)}
\NormalTok{summary\_table }\OtherTok{\textless{}{-}} \FunctionTok{data.frame}\NormalTok{(}
  \AttributeTok{scenario =} \FunctionTok{c}\NormalTok{(}\StringTok{\textquotesingle{}low\textquotesingle{}}\NormalTok{, }\StringTok{\textquotesingle{}high\textquotesingle{}}\NormalTok{),}
  \AttributeTok{R =} \FunctionTok{c}\NormalTok{(final\_low}\SpecialCharTok{$}\NormalTok{R, final\_high}\SpecialCharTok{$}\NormalTok{R),}
  \AttributeTok{F =} \FunctionTok{c}\NormalTok{(final\_low}\SpecialCharTok{$}\NormalTok{F, final\_high}\SpecialCharTok{$}\NormalTok{F),}
  \AttributeTok{A =} \FunctionTok{c}\NormalTok{(final\_low}\SpecialCharTok{$}\NormalTok{A, final\_high}\SpecialCharTok{$}\NormalTok{A),}
  \AttributeTok{m =} \FunctionTok{c}\NormalTok{(final\_low}\SpecialCharTok{$}\NormalTok{m, final\_high}\SpecialCharTok{$}\NormalTok{m)}
\NormalTok{)}

\NormalTok{knitr}\SpecialCharTok{::}\FunctionTok{kable}\NormalTok{(summary\_table, }\AttributeTok{digits =} \DecValTok{6}\NormalTok{)}
\end{Highlighting}
\end{Shaded}

\begin{longtable}[]{@{}lrrrr@{}}
\toprule\noalign{}
scenario & R & F & A & m \\
\midrule\noalign{}
\endhead
\bottomrule\noalign{}
\endlastfoot
low & 1.045455 & 99.64624 & 9.696967 & 0.550000 \\
high & 1.800536 & 98.59535 & 36.939699 & 0.383882 \\
\end{longtable}

\section{Results and Discussion}\label{results-and-discussion}

The R code above produces time series plots for fish and lamprey
populations under low and high resource regimes, and a male fraction
plot. It also prints a summary table with final values for R, F, A, and
m under each scenario.

\section{Appendix: Full code}\label{appendix-full-code}

The R code used for the model is included in the model code chunk above.
To render this file to PDF or HTML, knit as usual in RStudio or run.

\end{document}
